\section{Quantitative compactness through certified finite nets}
\label{sec:finite-nets}

Classical analysis often invokes compactness to pass from local information to a
uniform conclusion or to obtain an extremum.  In a finitary setting, compactness is an
interface consisting of a finite net and a modulus.  This section develops the generic
minimum theorem; the same pattern supports uniform continuity, covering arguments,
finite-dimensional topology, and integration.

\subsection{Net presentation}

Let \((X,d)\) be a represented metric object.  A certified net interface supplies, for
each precision parameter \(k\):

\begin{itemize}
\item a finite list \(N_k=(p_{k,0},\ldots,p_{k,m(k)-1})\);
\item a radius \(\delta(k)\in\Q_{>0}\); and
\item a finite certificate that every admissible input observation is within
\(\delta(k)\) of some net point.
\end{itemize}

For a function \(f:X\to\mathbb R\), a modulus supplies a rational error
\(\omega(k)\) such that
\[
  d(x,y)\le\delta(k)\quad\Longrightarrow\quad
  \abs{f(x)-f(y)}\le\omega(k).
\]
The evaluator returns rationals \(v_{k,i}\) with
\[
  \abs{v_{k,i}-f(p_{k,i})}\le\eta(k).
\]
All finite data and bounds are primitive recursive in \(k\).

\subsection{Finite minimum algorithm}

Choose the least index \(i_*\) minimizing \(v_{k,i}\).  This is a primitive-recursive
fold or bounded search.  Put \(m_k=v_{k,i_*}\) and \(p_k=p_{k,i_*}\).

\begin{theorem}[Finite-net minimum]
\label{thm:finite-net-min}
Under the net, modulus, and evaluator contracts:
\begin{enumerate}[label=(\roman*),leftmargin=2.4em]
\item for every admissible \(x\in X\),
\[
   f(p_k)\le f(x)+\omega(k)+2\eta(k);
\]
\item if \(\inf_X f\) is used in the external semantics, then
\[
   \inf_X f-\eta(k)\le m_k
   \le \inf_X f+\omega(k)+\eta(k);
\]
\item the selected point is an
\((\omega(k)+2\eta(k))\)-approximate minimizer; and
\item the finite comparison theorem, with \(x\) represented by an observation and its
coverage witness, is derivable in \(\iPRA\) using only bounded search and rational
arithmetic.
\end{enumerate}
\end{theorem}

\begin{proof}
Fix admissible \(x\).  By coverage, choose a net point \(p_{k,j}\) with
\(d(x,p_{k,j})\le\delta(k)\).  The modulus gives
\(f(p_{k,j})\le f(x)+\omega(k)\).  Evaluation error gives
\[
   v_{k,j}\le f(p_{k,j})+\eta(k)
            \le f(x)+\omega(k)+\eta(k).
\]
Minimality of \(i_*\) yields \(m_k\le v_{k,j}\), and another evaluation estimate gives
\[
   f(p_k)\le m_k+\eta(k)
           \le f(x)+\omega(k)+2\eta(k),
\]
proving (i) and (iii).

For the lower bound in (ii), every net value satisfies
\(v_{k,i}\ge f(p_{k,i})-\eta(k)\ge\inf_X f-\eta(k)\), so their minimum does too.  For
the upper bound, apply the preceding inequality to points whose function values
approach the infimum, or equivalently take the infimum over \(x\) in the bound for
\(m_k\).

For (iv), avoid quantifying over an abstract real space.  The input supplies a finite
observation of \(x\) and a bounded index witnessing net coverage.  Booleanize the
coverage and rational inequalities, instantiate the modulus contract, and use the
finite-minimum comparison proved by induction over the net list.  The resulting object
formula says that the selected sample is no worse, up to the displayed error, than any
input carrying a coverage certificate.  Its external interpretation is (i).
\end{proof}

\subsection{No exact argmin is claimed}

An approximate minimum value is stable under small perturbations; an exact minimizing
point need not be.  A flat interval, two equal wells, or a symmetric problem may have
many minimizers.  Therefore the generic interface returns a value name and approximate
minimizer cloud, not an exact point.

If a quantitative uniqueness modulus is available, coherent minimizers can be
constructed.  One useful form is
\[
   f(x)\le \inf f+\alpha(p)
   \ \wedge\ 
   f(y)\le \inf f+\alpha(p)
   \quad\Longrightarrow\quad d(x,y)\le2^{-p}.
\]
Choose net/error parameters making the approximation error at most \(\alpha(p)\).
Then all approximate minimizers lie in a ball of radius \(2^{-p}\), providing the
diameter bound required by shrinking viability.

\subsection{Strong convexity as a uniqueness producer}

For a midpoint strongly convex function with parameter \(\lambda>0\), suppose
\[
  f\!\left(\frac{x+y}{2}\right)
  \le\frac{f(x)+f(y)}2-\frac{\lambda}{8}d(x,y)^2.
\]
Since the midpoint value is at least \(\inf f\), if both \(x\) and \(y\) are within
\(\alpha\) of the infimum, then
\[
  \frac{\lambda}{8}d(x,y)^2\le\alpha,
  \qquad d(x,y)\le\sqrt{8\alpha/\lambda}.
\]
Choosing dyadic \(\alpha\) with an explicit rational square bound yields a
primitive-recursive uniqueness modulus.  Thus strong convexity legitimately upgrades
approximate minimizer clouds to a coherent point name.

\subsection{Other compactness templates}

The same interface yields:

\begin{itemize}
\item suprema and oscillation bounds by finite max/min folds;
\item uniform bounds from pointwise estimates plus a modulus;
\item Riemann sums and cubature from finite partitions and variation moduli;
\item Lebesgue-number-style finite refinements when an explicit cover margin is given;
\item finite simplicial approximations of compact metric presentations; and
\item finite subcover conclusions when a coverage certificate is part of the input.
\end{itemize}

It does not extract a modulus or finite net from a bare qualitative assertion of
compactness.  The net producer is mathematical data.
